\chapter{\texorpdfstring{$(P_5,\text{house})$}{(P5,house)}-free graphs with \texorpdfstring{$C_5$}{C5}}
\label{cap:c5}

\section{Characterization of monopolarity}
\label{sec:c5_characterization}

Having solved the case for chordal graphs, now we turn our attention to
non-chordal graphs --- such graphs contain an induced copy of a cycle of length
at least 4. Since we only consider $P_5$-free graphs, we know that each induced
cycle has length at most 5. We devote this chapter to the family of graphs with
an induced copy of $C_5$. Note that, for all graphs in the present chapter with
an induced copy of $C_5$, the labelling of the vertices of the 5-cycle is
determined by the drawings shown. Thus, vertices of a 5-cycle are labeled from
$a$ to $e$, starting with the rightmost vertex and following a counter-clockwise
order (see \Cref{fig:labels_c5}).

\begin{figure}[htb!]
        \centering
        \begin{tikzpicture}[every node/.style={vertex}, node
            distance=1cm, on grid, edge, label distance=1.5pt, font=\footnotesize]
            \node[label={above:$b$}] (0) at ({90}:1){};
            \node[label={left:$c$}] (1)  at ({(360/5) + 90}:1){};
            \node[label={below left:$d$}] (2)  at ({(360/5)*2 + 90}:1){};
            \node[label={below right:$e$}] (3)  at ({(360/5)*3 + 90}:1){};
            \node[label={right:$a$}] (4)  at ({(360/5)*4 + 90}:1){};

            \foreach \i in {0,...,4}
                \draw let \n1={int(mod(\i+1,5))} in (\i) to (\n1);
        \end{tikzpicture}
        \caption{Labels considered for $C_5$.}
        \label{fig:labels_c5}
    \end{figure}

Recall that the final algorithm first checks if the input graph is chordal. If
it finds an induced copy of $C_5$, then it executes the processes described in
the following section. Otherwise, if it detects an induced copy of $C_4$, then
it proceeds to execute the instructions that will be described in the next
chapter. If at any point of such instructions the algorithm finds an induced
copy of $C_5$ (that, by chance, did not detect the algorithm that verifies
chordality), then it switches to do the procedures of the present chapter.

We start our analysis with a useful observation.

\begin{proposition}
    \label{prop:partition_c5}
    Up to symmetry, there is only one monopolar partition of $C_5$.
\end{proposition}

\begin{proof}
    Let $(A,B)$ be a monopolar partition of $C= \p{a,b,c,d,e,a}$. Since $C$
    contains an induced copy of $P_3$, at least one vertex must be in the
    independent part, say $a\in A$. Thus, $b$ and $e$ are in $B$. Since
    $\p{b,c,d,e}$ is an induced copy of $P_4$ and $c$ and $d$ are adjacent,
    exactly one between $c$ and $d$ is in $A$ and the other is in $B$, say $c$
    is in $A$. Therefore, $(\l{a,c},\l{b,d,e})$ is a monopolar partition of $C$,
    and all partitions of it are constructed in this way.

\end{proof}

This result gives us insight for deciding when a graph $G$ is monopolar: if it
is monopolar, then we know how each induced copy of $C_5$ within it must be
partitioned (up to symmetry). Meanwhile, if $G$ is not monopolar, we can prove
it by exhibiting an induced copy of $C_5$ within $G$ such that none of the
possible monopolar partitions of $C_5$ can be extended to a monopolar partition
of $G$. This last technique is used in the following proposition to introduce
new minimal monopolar obstructions.

\begin{figure}[htb!]
    \centering
    \parbox{0.8\linewidth}{
        \begin{tikzpicture}[every node/.style={vertex}, node
            distance=1cm, on grid, edge, baseline=(current bounding box.center)]
            \foreach \i in {0,...,4}
                \node (\i)  at ({(360/5)*\i + 90}:1){};

            \node (5) at (0,0) {};
            \node (H) [Label] at (0,-1.5) {$H_8$};
                
                \foreach \i in {0,...,4}{
                    \draw let \n1={int(mod(\i+1,5))} in (\i) to (\n1);
                    \draw (\i) to (5);
                }
        \end{tikzpicture}
        \hfill
        \begin{tikzpicture}[every node/.style={vertex}, node
            distance=1cm, on grid, edge, baseline=(current bounding box.center)]
                \foreach \i in {0,...,4}
                    \node (\i)  at ({(360/5)*\i + 90}:1){};
                
                \node (5) at (0,0.5) {};
                \node (6) at ({(360/5)*3 + 90}:0.5) {};
                \node (7) at (0,0) {};
                \node (H) [Label] at (0,-1.5) {$H_9$};

                    
                \foreach \i in {0,...,4}
                    \draw let \n1={int(mod(\i+1,5))} in (\i) to (\n1);
                \draw (1) to (5);
                \draw (4) to (5);
                \draw (4) to (6);
                \draw (2) to (6);
                \draw (3) to (6);
                \draw (4) to (7);
                \draw (2) to (7);
        \end{tikzpicture}
        \hfill
        \begin{tikzpicture}[every node/.style={vertex}, node
            distance=1cm, on grid, edge, baseline=(current bounding box.center)]
                \foreach \i in {0,...,4}
                    \node (\i)  at ({(360/5)*\i + 90}:1){};
                \node (5) at (0,0.5) {};
                \node (6) at ({(360/5)*3 + 90}:0.5) {};
                \node (7) at ({(360/5)*2 + 90}:0.5) {};
                \node (H) [Label] at (0,-1.5) {$H_{10}$};
                    
                \foreach \i in {0,...,4}
                    \draw let \n1={int(mod(\i+1,5))} in (\i) to (\n1);
                \draw (1) to (5);
                \draw (4) to (5);
                \draw (4) to (6);
                \draw (2) to (6);
                \draw (1) to (7);
                \draw (3) to (7);
                \draw (6) to (7);
        \end{tikzpicture}
        \hfill
        
        \begin{tikzpicture}[every node/.style={vertex}, node
            distance=1cm, on grid, edge, baseline=(current bounding box.center)]
            \foreach \i in {0,...,3}
                \node (\i)  at ({(360/4)*\i + 90}:1){};
            \node (4) at (0,0) {};
            \node (H) [Label] at (0,-1.5) {$H_{11}$};
                
            \foreach \i in {0,...,3}{
                \draw let \n1={int(mod(\i+1,4))} in (\i) to (\n1);
                \draw (\i) to (4);
            }
        \end{tikzpicture}
        \hfill
        \begin{tikzpicture}[every node/.style={vertex}, node
            distance=1cm, on grid, edge, baseline=(current bounding box.center)]
                \node (1) {};
                \node (2) [right= of 1] {};
                \node (3) [right= of 2] {};
                \node (4) [right= of 3] {};
                \node (5) [below= of 1] {};
                \node (6) [below= of 2] {};
                \node (7) [below= of 3] {};
                \node (8) [below= of 4] {};
                \node (H) [Label] at (1.5,-1.5) {$H_{12}$};

                \draw (1) to (2);
                \draw (2) to (3);
                \draw (3) to (4);
                \draw (1) to (5);
                \draw (4) to (8);
                \draw (5) to (6);
                \draw (6) to (7);
                \draw (7) to (8);
                \draw (1) to (6);
                \draw (2) to (5);
                \draw (2) to (7);
                \draw (3) to (6);
                \draw (3) to (8);
                \draw (4) to (7);
        \end{tikzpicture}
        \hfill
        \begin{tikzpicture}[every node/.style={vertex}, node
            distance=1cm, on grid, edge, baseline=(current bounding box.center)]
                \node (1) {};
                \node (2) [right= of 1] {};
                \node (3) [below= of 2] {};
                \node (4) [left= of 3] {};
                \node (5) at ($(4) + (-150:1)$) {};
                \node (6) at ($(3) + (-30:1)$) {};
                \node (H) [Label] at (0.5,-2) {$H_{13}$};

                \draw (1) to (2);
                \draw (1) to (3);
                \draw (1) to (4);
                \draw (2) to (3);
                \draw (2) to (4);
                \draw (3) to (4);
                \draw (6) to (2);
                \draw (6) to (3);
                \draw (5) to (1);
                \draw (5) to (4);
                \draw (5) to (6);

        \end{tikzpicture}
    }
    
    \caption{Some minimal obstructions for monopolarity in
    $(P_5,\text{house})$-free graphs containing an induced copy of $C_5$.}
    \label{fig:p5-house-obstruction-with-c5}
\end{figure}

\begin{proposition}
    \label{prop:minimal_obs_c5}
    The graphs depicted in \Cref{fig:p5-house-obstruction-with-c5} are minimal
    monopolar obstructions.
\end{proposition}

\begin{proof}
    Suppose $H_8$ is monopolar. By \Cref{prop:partition_c5}, without loss of
    generality the monopolar partition induced on its induced copy of $C_5$ is
    $(\l{a,c},\l{b,d,e})$. Then, $u$, the universal vertex of $H_8$, cannot be
    in $A$ since $a$ is adjacent to $u$, and cannot be in $B$, since $\p{b,u,d}$
    is an induced copy of $P_3$. It follows that $H_8$ is not monopolar.
    Moreover, in \Cref{fig:minimal_h8} the monopolar partitions of the unique
    graphs obtained from deleting a vertex from $H_8$ (up to isomorphism) are
    shown.

    \begin{figure}[htb!]
        \centering
        \begin{tikzpicture}
            \matrix[column sep=0.5cm]{
                \node{
                    \begin{tikzpicture}[every node/.style={vertex}, node
                        distance=1cm, on grid, edge]
                        \foreach \i in {0,...,4}
                            \node (\i)  at ({(360/5)*\i + 90}:1){};
                        
                        \draw (0) to (1);
                        \draw (1) to (2);
                        \draw (2) to (3);
                        \draw (3) to (4);
                        \draw (0) to (4);
        
                        \begin{scope}[on background layer, scale=1.5]
                            \node [fit = (4), draw=none, fill=blue!20, inner sep =
                            0.5pt] {};
                            \node [fit = (1), draw=none, fill=blue!20, inner sep =
                            0.5pt] {};
                            \node [fit = (0), draw=none, fill=red!20, inner sep =
                            0.5pt] {};
                            \draw [draw=none, fill=red!20] \convexpath{2,3}{3.5pt};
                        \end{scope}
                    \end{tikzpicture}
                }; &
                \node{
                    \begin{tikzpicture}[every node/.style={vertex}, node
                        distance=1cm, on grid, edge]
                        \node[Label] (0) at (90:1) {};
                        \node (1) at (162:1) {};
                        \node (2) at (234:1) {};
                        \node (3) at (306:1) {};
                        \node (4) at (18:1) {};
                        \node (5) {};
                        
                        \draw (1) to (2);
                        \draw (2) to (3);
                        \draw (3) to (4);
                        \draw (1) to (5);
                        \draw (2) to (5);
                        \draw (3) to (5);
                        \draw (4) to (5);
        
                        \begin{scope}[on background layer, scale=1.5]
                            \node [fit = (4), draw=none, fill=blue!20, inner sep =
                            0.5pt] {};
                            \node [fit = (1), draw=none, fill=blue!20, inner sep =
                            0.5pt] {};
                            \draw [draw=none, fill=red!20] \convexpath{2,5,3}{3.5pt};
                        \end{scope}
                    \end{tikzpicture}
                }; \\
            };
        \end{tikzpicture}
        \caption{Monopolar partitions of $H_8 - v$.}
        \label{fig:minimal_h8}
    \end{figure}

    Consider the following labelling of the vertices of $H_9$ (see
    \Cref{fig:labelling_h9}).
    \begin{figure}[htb!]
        \centering
        \begin{tikzpicture}[every node/.style={vertex}, node
            distance=1cm, on grid, edge, label distance=1.5pt, font=\footnotesize]
            \node[label={above:$b$}] (0) at ({90}:1){};
            \node[label={left:$c$}] (1)  at ({(360/5) + 90}:1){};
            \node[label={below left:$d$}] (2)  at ({(360/5)*2 + 90}:1){};
            \node[label={below right:$e$}] (3)  at ({(360/5)*3 + 90}:1){};
            \node[label={right:$a$}] (4)  at ({(360/5)*4 + 90}:1){};
            \node[label={[left=0pt]90:$b'$}] (5) at (0,0.5) {};
            \node[label={above:$e'$}] (6) at ({(360/5)*3 + 90}:0.5) {};
            \node[label={left:$e''$}] (7) at (0,0) {};

            \foreach \i in {0,...,4}
                \draw let \n1={int(mod(\i+1,5))} in (\i) to (\n1);
            \draw (1) to (5);
            \draw (4) to (5);
            \draw (4) to (6);
            \draw (2) to (6);
            \draw (3) to (6);
            \draw (4) to (7);
            \draw (2) to (7);
        \end{tikzpicture}
        \caption{A labelling of $H_9$.}
        \label{fig:labelling_h9}
    \end{figure}
    Suppose $H_9$ is monopolar. Notice $\l{a,d,e,e'}$ induces a copy of the
    diamond. Therefore, \Cref{lem:monopolar_diamond} implies that both $e$ and
    $e'$ must be in $B$. It follows that at least one between $a$ and $d$ must
    be in $A$, otherwise, $B$ would contain an induced copy of $P_3$. Since $A$
    is stable, it follows that $e''$ must be in $B$. Meanwhile, $G[B]$ being a
    cluster implies that both $a$ and $d$ are in $A$. Since $A$ is stable and
    $b,b',c\in N(a)\cup N(d)$, it follows that $\p{b,c,b'}$ is an induced copy
    of $P_3$ in $B$, once more contradicting that $G[B]$ is a cluster. Then,
    $H_9$ is not monopolar, and in \Cref{fig:minimal_h9} the monopolar
    partitions of $H_9 - v$ are shown.

    \begin{figure}[htb!]
        \centering
        \begin{tikzpicture}
            \matrix[column sep=0.5cm, row sep=0.2cm]{
                \node{
                    \begin{tikzpicture}[every node/.style={vertex}, node
                        distance=1cm, on grid, edge]
                        \node (0) at (90:1) {};
                        \node (1)  at ({(360/5) + 90}:1){};
                        \node (2)  at ({(360/5)*2 + 90}:1){};
                        \node (3)  at ({(360/5)*3 + 90}:1){};
                        \node (4)  at ({(360/5)*4 + 90}:1){};
                        \node (6) at ({(360/5)*3 + 90}:0.5) {};
                        \node (7) at (0,0) {};
            
                        \draw (0) to (1);
                        \draw (1) to (2);
                        \draw (2) to (3);
                        \draw (3) to (4);
                        \draw (0) to (4);
                        \draw (4) to (6);
                        \draw (2) to (6);
                        \draw (3) to (6);
                        \draw (4) to (7);
                        \draw (2) to (7);
            
                        \begin{scope}[on background layer, scale=1.5]
                            \draw [draw=none, fill=red!20] \convexpath{3,6}{3.5pt};
                            \draw [draw=none, fill=red!20] \convexpath{0,1}{3.5pt};
                            \node [fit=(7), draw=none, fill=red!20, inner sep=0.5pt] {};
                            \node [fit=(2), draw=none, fill=blue!20, inner sep=0.5pt] {};
                            \node [fit=(4), draw=none, fill=blue!20, inner sep=0.5pt] {};
                        \end{scope}
                    \end{tikzpicture}
                }; &
                \node{
                    \begin{tikzpicture}[every node/.style={vertex}, node
                        distance=1cm, on grid, edge]
                        \node (0) at ({90}:1){};
                        \node[Label] (1)  at ({(360/5) + 90}:1){};
                        \node (2)  at ({(360/5)*2 + 90}:1){};
                        \node (3)  at ({(360/5)*3 + 90}:1){};
                        \node (4)  at ({(360/5)*4 + 90}:1){};
                        \node (5) at (0,0.5) {};
                        \node (6) at ({(360/5)*3 + 90}:0.5) {};
                        \node (7) at (0,0) {};
            
                        \draw (2) to (3);
                        \draw (3) to (4);
                        \draw (0) to (4);
                        \draw (4) to (5);
                        \draw (4) to (6);
                        \draw (2) to (6);
                        \draw (3) to (6);
                        \draw (4) to (7);
                        \draw (2) to (7);
            
                        \begin{scope}[on background layer, scale=1.5]
                            \draw [draw=none, fill=red!20] \convexpath{3,6}{3.5pt};
                            \node [fit=(7), draw=none, fill=red!20, inner sep=0.5pt] {};
                            \node [fit=(0), draw=none, fill=red!20, inner sep=0.5pt] {};
                            \node [fit=(5), draw=none, fill=red!20, inner sep=0.5pt] {};
                            \node [fit=(2), draw=none, fill=blue!20, inner sep=0.5pt] {};
                            \node [fit=(4), draw=none, fill=blue!20, inner sep=0.5pt] {};
                        \end{scope}
                    \end{tikzpicture}
                }; &
                \node{
                    \begin{tikzpicture}[every node/.style={vertex}, node
                        distance=1cm, on grid, edge]
                        \node (0) at ({90}:1){};
                        \node (1)  at ({(360/5) + 90}:1){};
                        \node (2)  at ({(360/5)*2 + 90}:1){};
                        \node (3)  at ({(360/5)*3 + 90}:1){};
                        \node[Label] (4)  at ({(360/5)*4 + 90}:1){};
                        \node (5) at (0,0.5) {};
                        \node (6) at ({(360/5)*3 + 90}:0.5) {};
                        \node (7) at (0,0) {};
            
                        \draw (0) to (1);
                        \draw (1) to (2);
                        \draw (2) to (3);
                        \draw (1) to (5);
                        \draw (2) to (6);
                        \draw (3) to (6);
                        \draw (2) to (7);
            
                        \begin{scope}[on background layer, scale=1.5]
                            \draw [draw=none, fill=red!20] \convexpath{3,6}{3.5pt};
                            \node [fit=(7), draw=none, fill=red!20, inner sep=0.5pt] {};
                            \node [fit=(1), draw=none, fill=red!20, inner sep=0.5pt] {};
                            \node [fit=(2), draw=none, fill=blue!20, inner sep=0.5pt] {};
                            \node [fit=(0), draw=none, fill=blue!20, inner sep=0.5pt] {};
                            \node [fit=(5), draw=none, fill=blue!20, inner sep=0.5pt] {};
                        \end{scope}
                    \end{tikzpicture}
                }; \\
                \node{
                    \begin{tikzpicture}[every node/.style={vertex}, node
                        distance=1cm, on grid, edge]
                        \node (0) at ({90}:1){};
                        \node (1)  at ({(360/5) + 90}:1){};
                        \node (2)  at ({(360/5)*2 + 90}:1){};
                        \node (3)  at ({(360/5)*3 + 90}:1){};
                        \node (4)  at ({(360/5)*4 + 90}:1){};
                        \node (5) at (0,0.5) {};
                        \node (6) at ({(360/5)*3 + 90}:0.5) {};
                        \node[Label] (7) at (0,0) {};
            
                        \draw (0) to (1);
                        \draw (1) to (2);
                        \draw (2) to (3);
                        \draw (3) to (4);
                        \draw (0) to (4);
                        \draw (1) to (5);
                        \draw (4) to (5);
                        \draw (4) to (6);
                        \draw (2) to (6);
                        \draw (3) to (6);
            
                        \begin{scope}[on background layer, scale=1.5]
                            \draw [draw=none, fill=red!20] \convexpath{3,6,4}{3.5pt};
                            \node [fit=(1), draw=none, fill=red!20, inner sep=0.5pt] {};
                            \node [fit=(2), draw=none, fill=blue!20, inner sep=0.5pt] {};
                            \node [fit=(0), draw=none, fill=blue!20, inner sep=0.5pt] {};
                            \node [fit=(5), draw=none, fill=blue!20, inner sep=0.5pt] {};
                        \end{scope}
                    \end{tikzpicture}
                }; &
                \node{
                    \begin{tikzpicture}[every node/.style={vertex}, node
                        distance=1cm, on grid, edge]
                        \node (0) at ({90}:1){};
                        \node (1)  at ({(360/5) + 90}:1){};
                        \node (2)  at ({(360/5)*2 + 90}:1){};
                        \node (3)  at ({(360/5)*3 + 90}:1){};
                        \node (4)  at ({(360/5)*4 + 90}:1){};
                        \node (5) at (0,0.5) {};
                        \node[Label] (6) at ({(360/5)*3 + 90}:0.5) {};
                        \node (7) at (0,0) {};
            
                        \draw (0) to (1);
                        \draw (1) to (2);
                        \draw (2) to (3);
                        \draw (3) to (4);
                        \draw (0) to (4);
                        \draw (1) to (5);
                        \draw (4) to (5);
                        \draw (2) to (7);
                        \draw (4) to (7);
            
                        \begin{scope}[on background layer, scale=1.5]
                            \draw [draw=none, fill=red!20] \convexpath{1,2}{3.5pt};
                            \node [fit=(4), draw=none, fill=red!20, inner sep=0.5pt] {};
                            \node [fit=(3), draw=none, fill=blue!20, inner sep=0.5pt] {};
                            \node [fit=(0), draw=none, fill=blue!20, inner sep=0.5pt] {};
                            \node [fit=(5), draw=none, fill=blue!20, inner sep=0.5pt] {};
                            \node [fit=(7), draw=none, fill=blue!20, inner sep=0.5pt] {};
                        \end{scope}
                    \end{tikzpicture}
                }; &
                \node{
                    \begin{tikzpicture}[every node/.style={vertex}, node
                        distance=1cm, on grid, edge]
                        \node (0) at ({90}:1){};
                        \node (1)  at ({(360/5) + 90}:1){};
                        \node[Label] (2)  at ({(360/5)*2 + 90}:1){};
                        \node (3)  at ({(360/5)*3 + 90}:1){};
                        \node (4)  at ({(360/5)*4 + 90}:1){};
                        \node (5) at (0,0.5) {};
                        \node (6) at ({(360/5)*3 + 90}:0.5) {};
                        \node (7) at (0,0) {};
            
                        \draw (0) to (1);
                        \draw (3) to (4);
                        \draw (0) to (4);
                        \draw (1) to (5);
                        \draw (4) to (5);
                        \draw (4) to (6);
                        \draw (3) to (6);
                        \draw (4) to (7);
            
                        \begin{scope}[on background layer, scale=1.5]
                            \draw [draw=none, fill=red!20] \convexpath{3,6,4}{3.5pt};
                            \node [fit=(7), draw=none, fill=blue!20, inner sep=0.5pt] {};
                            \node [fit=(0), draw=none, fill=blue!20, inner sep=0.5pt] {};
                            \node [fit=(5), draw=none, fill=blue!20, inner sep=0.5pt] {};
                            \node [fit=(1), draw=none, fill=red!20, inner sep=0.5pt] {};
                        \end{scope}
                    \end{tikzpicture}
                }; \\
            };
        \end{tikzpicture}
    \caption{Monopolar partitions of $H_9-v$.}
    \label{fig:minimal_h9}
    \end{figure}

    Consider the following labelling of $H_{10}$ (see \Cref{fig:labelling_h10}).
    \begin{figure}[htb!]
        \centering
        \begin{tikzpicture}[every node/.style={vertex}, node
            distance=1cm, on grid, edge, label distance=1.5pt, font=\footnotesize]
            \node[label={above:$b$}] (0) at ({90}:1){};
            \node[label={left:$c$}] (1)  at ({(360/5) + 90}:1){};
            \node[label={below left:$d$}] (2)  at ({(360/5)*2 + 90}:1){};
            \node[label={below right:$e$}] (3)  at ({(360/5)*3 + 90}:1){};
            \node[label={right:$a$}] (4)  at ({(360/5)*4 + 90}:1){};
            \node[label={[left=0pt]90:$b'$}] (5) at (0,0.5) {};
            \node[label={above:$e'$}] (6) at ({(360/5)*3 + 90}:0.5) {};
            \node[label={above:$d'$}] (7) at ({(360/5)*2 + 90}:0.5) {};
                
            \foreach \i in {0,...,4}
                \draw let \n1={int(mod(\i+1,5))} in (\i) to (\n1);
            \draw (1) to (5);
            \draw (4) to (5);
            \draw (4) to (6);
            \draw (2) to (6);
            \draw (1) to (7);
            \draw (3) to (7);
            \draw (6) to (7);
        \end{tikzpicture}
        \caption{A labelling of $H_{10}$.}
        \label{fig:labelling_h10}
    \end{figure}
    Suppose $H_{10}$ is monopolar and $d$ is in $A$. We know from
    \Cref{prop:partition_c5} that either $a$ or $b$ is in $A$. If $b\in A$,
    then, since $A$ is stable and $a,e,a'\in N(d) \cup N(b)$, it follows that
    $\p{e,a,e'}$ is an induced copy of $P_3$ in $B$. Analogously, if $a\in A$,
    then $\p{b,c,b'}$ is an induced copy of $P_3$ in $B$. Thus, $d$ cannot be in
    $A$, so $d\in B$. Similarly follows that $e\in B$. Then, for $G[B]$ to be
    $P_3$-free, both $d'$ and $e'$ must be in $A$, contradicting that $A$ is
    stable. In \Cref{fig:minimal_h10} the monopolar partitions of $H_{10}-v$ are
    shown.
    
    \begin{figure}[htb!]
        \centering
        \begin{tikzpicture}
            \matrix[column sep = 0.5cm]{
                \node{
                    \begin{tikzpicture}[every node/.style={vertex}, node
                        distance=1cm, on grid, edge]
                        \node (0) at ({90}:1){};
                        \node (1)  at ({(360/5) + 90}:1){};
                        \node (2)  at ({(360/5)*2 + 90}:1){};
                        \node (3)  at ({(360/5)*3 + 90}:1){};
                        \node (4)  at ({(360/5)*4 + 90}:1){};
                        \node[Label] (5) at (0,0.5) {};
                        \node (6) at ({(360/5)*3 + 90}:0.5) {};
                        \node (7) at ({(360/5)*2 + 90}:0.5) {};
                            
                        \draw (0) to (1);
                        \draw (1) to (2);
                        \draw (2) to (3);
                        \draw (3) to (4);
                        \draw (0) to (4);
                        \draw (4) to (6);
                        \draw (2) to (6);
                        \draw (1) to (7);
                        \draw (3) to (7);
                        \draw (6) to (7);
            
                        \begin{scope}[on background layer, scale=1.5]
                            \draw [draw=none, fill=red!20] \convexpath{0,1}{3.5pt};
                            \node [fit=(3), fill=red!20, inner sep=0.5pt, draw=none] {};
                            \node [fit=(6), fill=red!20, inner sep=0.5pt, draw=none] {};
                            \node [fit=(4), fill=blue!20, inner sep=0.5pt, draw=none] {};
                            \node [fit=(2), fill=blue!20, inner sep=0.5pt, draw=none] {};
                            \node [fit=(7), fill=blue!20, inner sep=0.5pt, draw=none] {};
                        \end{scope}
                    \end{tikzpicture}
                }; &
                \node{
                    \begin{tikzpicture}[every node/.style={vertex}, node
                        distance=1cm, on grid, edge]
                        \node (0) at ({90}:1){};
                        \node[Label] (1)  at ({(360/5) + 90}:1){};
                        \node (2)  at ({(360/5)*2 + 90}:1){};
                        \node (3)  at ({(360/5)*3 + 90}:1){};
                        \node (4)  at ({(360/5)*4 + 90}:1){};
                        \node (5) at (0,0.5) {};
                        \node (6) at ({(360/5)*3 + 90}:0.5) {};
                        \node (7) at ({(360/5)*2 + 90}:0.5) {};
            
                        \draw (2) to (3);
                        \draw (3) to (4);
                        \draw (0) to (4);
                        \draw (4) to (5);
                        \draw (4) to (6);
                        \draw (2) to (6);
                        \draw (3) to (7);
                        \draw (6) to (7);
            
                        \begin{scope}[on background layer, scale=1.5]
                            \node [fit=(0), fill=blue!20, inner sep=0.5pt, draw=none] {};
                            \node [fit=(3), fill=blue!20, inner sep=0.5pt, draw=none] {};
                            \node [fit=(5), fill=blue!20, inner sep=0.5pt, draw=none] {};
                            \node [fit=(6), fill=blue!20, inner sep=0.5pt, draw=none] {};
                            \node [fit=(4), fill=red!20, inner sep=0.5pt, draw=none] {};
                            \node [fit=(2), fill=red!20, inner sep=0.5pt, draw=none] {};
                            \node [fit=(7), fill=red!20, inner sep=0.5pt, draw=none] {};
                        \end{scope}
                    \end{tikzpicture}
                }; &
                \node{
                    \begin{tikzpicture}[every node/.style={vertex}, node
                        distance=1cm, on grid, edge]
                        \node (0) at ({90}:1){};
                        \node (1)  at ({(360/5) + 90}:1){};
                        \node (2)  at ({(360/5)*2 + 90}:1){};
                        \node (3)  at ({(360/5)*3 + 90}:1){};
                        \node (4)  at ({(360/5)*4 + 90}:1){};
                        \node (5) at (0,0.5) {};
                        \node (6) at ({(360/5)*3 + 90}:0.5) {};
                        \node[Label] (7) at ({(360/5)*2 + 90}:0.5) {};
                            
                        \draw (0) to (1);
                        \draw (1) to (2);
                        \draw (2) to (3);
                        \draw (3) to (4);
                        \draw (0) to (4);
                        \draw (1) to (5);
                        \draw (4) to (5);
                        \draw (4) to (6);
                        \draw (2) to (6);
            
                        \begin{scope}[on background layer, scale=1.5]
                            \draw [draw=none, fill=red!20] \convexpath{1,2}{3.5pt};
                            \node [fit=(0), fill=blue!20, inner sep=0.5pt, draw=none] {};
                            \node [fit=(3), fill=blue!20, inner sep=0.5pt, draw=none] {};
                            \node [fit=(5), fill=blue!20, inner sep=0.5pt, draw=none] {};
                            \node [fit=(6), fill=blue!20, inner sep=0.5pt, draw=none] {};
                            \node [fit=(4), fill=red!20, inner sep=0.5pt, draw=none] {};
                        \end{scope}
                    \end{tikzpicture}
                }; \\
                };
            \end{tikzpicture}
        \caption{Monopolar partitions of $H_{10}-v$.}
        \label{fig:minimal_h10}
    \end{figure}

    Suppose that $H_{11}$ is monopolar. Similar to the case of $H_{8}$, $u$, the
    universal vertex of $H_{11}$, must be in $B$ and at least one vertex
    different from $u$ must be in $A$, since $C_4$ is not a cluster. Suppose
    $b\in A$, so $\p{a,u,c}$ is an induced copy of $P_3$ in $B$. Therefore,
    $H_{11}$ is not monopolar. In \Cref{fig:minimal_h11} the monopolar
    partitions of $H_{11}-v$ are shown.

    \begin{figure}[htb!]
        \centering
        \begin{tikzpicture}
            \matrix [column sep = 0.5cm]{
                \node{
                    \begin{tikzpicture}[every node/.style={vertex}, node
                        distance=1cm, on grid, edge]
                        \foreach \i in {0,...,3}
                            \node (\i)  at ({(360/4)*\i + 90}:1){};
                        
                        \draw (0) to (1);
                        \draw (1) to (2);
                        \draw (2) to (3);
                        \draw (0) to (3);
                            
                        \begin{scope}[on background layer, scale=1.5]
                            \node [fit=(0), draw=none, fill=blue!20, inner sep=0.5pt] {};
                            \node [fit=(2), draw=none, fill=blue!20, inner sep=0.5pt] {};
                            \node [fit=(1), draw=none, fill=red!20, inner sep=0.5pt] {};
                            \node [fit=(3), draw=none, fill=red!20, inner sep=0.5pt] {};
                        \end{scope}
                    \end{tikzpicture}

                }; &
                \node{
                    \begin{tikzpicture}[every node/.style={vertex}, node
                        distance=1cm, on grid, edge]
                        \node[Label] (0) at (90:1) {};
                        \node (1) at (180:1) {};
                        \node (2) at (270:1) {};
                        \node (3) at (0:1) {};
                        \node (4) {};
                        
                        \draw (1) to (2);
                        \draw (2) to (3);
                        \draw (1) to (4);
                        \draw (2) to (4);
                        \draw (3) to (4);
                            
                        \begin{scope}[on background layer, scale=1.5]
                            \draw [draw=none, fill=red!20] \convexpath{2,4}{3.5pt};
                            \node [fit=(1), draw=none, fill=blue!20, inner sep=0.5pt] {};
                            \node [fit=(3), draw=none, fill=blue!20, inner sep=0.5pt] {};
                        \end{scope}
                    \end{tikzpicture}
                }; \\
            };
        \end{tikzpicture}
        \caption{Monopolar partitions of $H_{11}$.}
        \label{fig:minimal_h11}
    \end{figure}

    Consider the following labelling of $H_{12}$ (see \Cref{fig:labelling_h12}).
    \begin{figure}[htb!]
        \centering
        \begin{tikzpicture}[every node/.style={vertex}, node
            distance=1cm, on grid, edge, label distance=1.5pt, font=\footnotesize]
            \node (1) [label={above:$a$}] {};
            \node (2) [right= of 1,label={above:$b$}] {};
            \node (3) [right= of 2,label={above:$c$}] {};
            \node (4) [right= of 3,label={above:$d$}] {};
            \node (5) [below= of 1,label={below:$e$}] {};
            \node (6) [below= of 2,label={below:$f$}] {};
            \node (7) [below= of 3,label={below:$g$}] {};
            \node (8) [below= of 4,label={below:$h$}] {};

            \draw (1) to (2);
            \draw (2) to (3);
            \draw (3) to (4);
            \draw (1) to (5);
            \draw (4) to (8);
            \draw (5) to (6);
            \draw (6) to (7);
            \draw (7) to (8);
            \draw (1) to (6);
            \draw (2) to (5);
            \draw (2) to (7);
            \draw (3) to (6);
            \draw (3) to (8);
            \draw (4) to (7);
        \end{tikzpicture}
        \caption{A labelling of $H_{12}$.}
        \label{fig:labelling_h12}
    \end{figure}
    Suppose $H_{12}$ is monopolar. Notice $a$ must be in $B$, otherwise, $a$
    being in $A$ and $A$ being stable implies that $\p{b,e,f}$ is an induced
    copy of $P_3$ in $B$. Additionally, at least one vertex between $b$ and $f$
    must be in $A$. Similarly, $d$ must be in $B$ and at least one vertex
    between $c$ and $g$ must be in $A$. Since $b$ and $f$ are both adjacent to
    $c$ and $g$, it follows that $A$ is not a stable set. Hence, $H_{12}$ is not
    monopolar. In \Cref{fig:minimal_h12} the monopolar partitions of $H_{12} -
    v$ are shown.

    \begin{figure}[htb!]
        \centering
        \begin{tikzpicture}
            \matrix[column sep = 0.5cm]{
                \node{
                    \begin{tikzpicture}[every node/.style={vertex}, node 
                        distance=1cm, on grid, edge]
                        \node[Label] (1) {};
                        \node (2) [right= of 1] {};
                        \node (3) [right= of 2] {};
                        \node (4) [right= of 3] {};
                        \node (5) [below= of 1] {};
                        \node (6) [below= of 2] {};
                        \node (7) [below= of 3] {};
                        \node (8) [below= of 4] {};
            
                        \draw (2) to (3);
                        \draw (3) to (4);
                        \draw (4) to (8);
                        \draw (5) to (6);
                        \draw (6) to (7);
                        \draw (7) to (8);
                        \draw (2) to (5);
                        \draw (2) to (7);
                        \draw (3) to (6);
                        \draw (3) to (8);
                        \draw (4) to (7);
            
                        \begin{scope}[on background layer, scale=1.5]
                            \node [fit=(5), draw=none, fill=blue!20, inner sep=0.5pt] {};
                            \node [fit=(2), draw=none, fill=red!20, inner sep=0.5pt] {};
                            \node [fit=(6), draw=none, fill=red!20, inner sep=0.5pt] {};
                            \node [fit=(3), draw=none, fill=blue!20, inner sep=0.5pt] {};
                            \node [fit=(7), draw=none, fill=blue!20, inner sep=0.5pt] {};
                            \draw [draw=none, fill=red!20] \convexpath{4,8}{3.5pt};
                        \end{scope}
                    \end{tikzpicture}
                }; &
                \node{
                    \begin{tikzpicture}[every node/.style={vertex}, node
                        distance=1cm, on grid, edge]
                        \node (1) {};
                        \node (2) [right= of 1, Label] {};
                        \node (3) [right= of 2] {};
                        \node (4) [right= of 3] {};
                        \node (5) [below= of 1] {};
                        \node (6) [below= of 2] {};
                        \node (7) [below= of 3] {};
                        \node (8) [below= of 4] {};
            
                        \draw (3) to (4);
                        \draw (1) to (5);
                        \draw (4) to (8);
                        \draw (5) to (6);
                        \draw (6) to (7);
                        \draw (7) to (8);
                        \draw (1) to (6);
                        \draw (3) to (6);
                        \draw (3) to (8);
                        \draw (4) to (7);
            
                        \begin{scope}[on background layer, scale=1.5]
                            \node [fit=(5), draw=none, fill=blue!20, inner sep=0.5pt] {};
                            \node [fit=(3), draw=none, fill=blue!20, inner sep=0.5pt] {};
                            \node [fit=(7), draw=none, fill=blue!20, inner sep=0.5pt] {};
                            \draw [draw=none, fill=red!20] \convexpath{4,8}{3.5pt};
                            \draw [draw=none, fill=red!20] \convexpath{1,6}{3.5pt};
                        \end{scope}
                    \end{tikzpicture}
                }; \\
            };
        \end{tikzpicture}
        \caption{Monopolar partitions of $H_{12} - v$.}
        \label{fig:minimal_h12}
    \end{figure}

    Suppose $H_{13}$ is monopolar. Notice every vertex of the induced copy of
    $K_4$ must be in $B$, for if any was in $A$, then two vertices of the
    induced copy of $K_4$ and one not in there would create an induced copy of
    $P_3$ in $B$. Then, for $G[B]$ to be $P_3$-free, the two vertices no in the
    induced copy of $K_4$ must be in $A$, contradicting that $A$ is stable.
    Thus, $H_{13}$ is not monopolar. In \Cref{fig:minimal_h13} the monopolar
    partitions of $H_{13} - v$ are shown.

    \begin{figure}[htb!]
        \centering
        \begin{tikzpicture}
            \matrix[column sep=0.5cm]{
                \node{
                    \begin{tikzpicture}[every node/.style={vertex}, node
                        distance=1cm, on grid, edge]
                        \node[Label] (1) {};
                        \node (2) [right= of 1] {};
                        \node (3) [below= of 2] {};
                        \node (4) [left= of 3] {};
                        \node (5) at ($(4) + (-150:1)$) {};
                        \node (6) at ($(3) + (-30:1)$) {};
            
                        \draw (2) to (3);
                        \draw (2) to (4);
                        \draw (3) to (4);
                        \draw (6) to (2);
                        \draw (6) to (3);
                        \draw (5) to (4);
                        \draw (5) to (6);
            
                        \begin{scope}[on background layer, scale=1.5]
                            \draw [draw=none, fill=red!20] \convexpath{2,3}{3.5pt};
                            \node [fit=(5), draw=none, fill=red!20, inner sep=0.5pt] {};
                            \node [fit=(4), draw=none, fill=blue!20, inner sep=0.5pt] {};
                            \node [fit=(6), draw=none, fill=blue!20, inner sep=0.5pt] {};
                        \end{scope}
                    \end{tikzpicture}
                }; &
                \node{
                    \begin{tikzpicture}[every node/.style={vertex}, node
                        distance=1cm, on grid, edge]
                        \node (1) {};
                        \node (2) [right= of 1] {};
                        \node (3) [below= of 2] {};
                        \node (4) [left= of 3] {};
                        \node[Label] (5) at ($(4) + (-150:1)$) {};
                        \node (6) at ($(3) + (-30:1)$) {};
            
                        \draw (1) to (2);
                        \draw (1) to (3);
                        \draw (1) to (4);
                        \draw (2) to (3);
                        \draw (2) to (4);
                        \draw (3) to (4);
                        \draw (6) to (2);
                        \draw (6) to (3);
            
                        \begin{scope}[on background layer, scale=1.5]
                            \draw [draw=none, fill=red!20] \convexpath{1,2,3,4}{3.5pt};
                            \node [fit=(6), draw=none, fill=blue!20, inner sep=0.5pt] {};
                        \end{scope}
                    \end{tikzpicture}
                }; \\
            };
        \end{tikzpicture}
        \caption{Monopolar partitions of $H_{13}-v$.}
        \label{fig:minimal_h13}
    \end{figure}
\end{proof}

We will prove that graphs $H_1$ all the way to $H_{13}$ characterize
monopolarity for the family of $(P_5,\text{house})$-free graphs with an induced
copy of $C_5$. To do so, we prove a series of lemmas describing the exact
neighborhoods of a 5-cycle. Throughout the following proofs we will always
denote an induced copy of $C_5$ in a graph $G$ by $C = \p{a,b,c,d,e,a}$ and the
set $\l{a,b,c,d,e}$ will be denoted by $Z$.

\begin{lemma}
    \label{lem:pre_neighborhoods_c5}
    Let $G$ be a $(P_5,\text{house})$-free graph with an induced copy of $C_5$.
    Let $X$ be a proper subset of $Z$.
    \begin{lemenum}
        \item If $X$ is nonempty and consists of at most two adjacent vertices
            of $C$, then $N_X$ is empty. \label{lem:pre_c5_1}
        \item If $X$ has three non-consecutive vertices, then $N_X$ is empty.
            \label{lem:pre_c5_2}
        \item If $X$ is nonempty, then $N_\varnothing$ is anticomplete to $N_X$.
            \label{lem:pre_c5_3}
    \end{lemenum}
\end{lemma}

\begin{proof}
    \leavevmode \par
    \begin{enumerate}[label=(\roman*)]
        \item Since $X$ is not empty, suppose without loss of generality that
            $a\in X$. Since $X$ does not contain nonadjacent vertices of $C$, it
            follows that $X$ has at most one additional vertex, either $b$ or
            $e$, but not both. Suppose without loss of generality that
            $X\subseteq \l{a,e}$. If $x\in N_X$, then $\p{x,a,b,c,d}$ is an
            induced copy of $P_5$ (see \Cref{fig:pre_1}). Therefore, $N_X$ is
            empty.

            \begin{figure}[htb!]
                \centering
                \begin{tikzpicture}[every node/.style={vertex}, node
                distance=1cm, on grid, edge, baseline=(current bounding box.center)]
                \foreach \i in {0,...,2,4}
                    \node (\i)  at ({(360/5)*\i + 90}:1){};
                \node[lightgray] (3)  at ({(360/5)*3 + 90}:1){};
                \node (5) at (0,0) {};
                    
                \draw[lightgray] (4) to (3);
                \draw[lightgray] (2) to (3);
                \draw[lightgray] (5) to (3);
                \draw (0) to (1);
                \draw (2) to (1);
                \draw (0) to (4);
                \draw (5) to (4);
                \end{tikzpicture}
                \caption{An induced copy of $P_5$ given that $N_X$ is nonempty.}
                \label{fig:pre_1}
            \end{figure}

        \item Since $X$ is not $Z$, suppose without loss of generality that $X$
            has $a$, $c$, $d$, and possibly $e$, but $b\notin X$. If $x\in N_X$,
            then $\l{a,b,c,x,d}$ induces a copy of the house (see
            \Cref{fig:pre_2}). Hence, $N_X$ is empty.

            \begin{figure}[htb!]
                \centering
                \begin{tikzpicture}[every node/.style={vertex}, node
                distance=1cm, on grid, edge, baseline=(current bounding box.center)]
                \foreach \i in {0,...,2,4}
                    \node (\i)  at ({(360/5)*\i + 90}:1){};
                \node[lightgray] (3)  at ({(360/5)*3 + 90}:1){};
                \node (5) at (0,0) {};
                    
                \draw[lightgray] (4) to (3);
                \draw[lightgray] (2) to (3);
                \draw[lightgray] (5) to (3);
                \draw (0) to (1);
                \draw (2) to (1);
                \draw (0) to (4);
                \draw (5) to (4);
                \draw (5) to (1);
                \draw (5) to (2);
                \end{tikzpicture}
                \caption{An induced copy of the house given that $N_X$ is nonempty.}
                \label{fig:pre_2}
            \end{figure}
        
        \item Since $X \neq Z$, by the previous two items, $X$ either consists
            of two nonadjacent vertices or three consecutive vertices. Thus,
            suppose $X\subseteq \l{a,b,c}$. If $w \in N_{\varnothing}$ is
            adjacent to $x\in N_X$, then $\p{w,x,a,e,d}$ is an induced copy of
            $P_5$ (see \Cref{fig:pre_3}). It follows that $N_{\varnothing}$ is
            anticomplete to $N_X$.

            \begin{figure}[htb!]
                \centering
                \begin{tikzpicture}[every node/.style={vertex}, node
                distance=1cm, on grid, edge]
                \foreach \i in {2,3,4}
                    \node (\i)  at ({(360/5)*\i + 90}:1){};
                \node[lightgray] (0)  at ({90}:1){};
                \node[lightgray] (1)  at ({(360/5) +90}:1){};
                \node (5) at (0,0.5) {};
                \node (6) at (0,0) {};
                    
                \draw[lightgray] (5) to (1);
                \draw[lightgray] (5) to (0);
                \draw[lightgray] (0) to (1);
                \draw[lightgray] (0) to (4);
                \draw[lightgray] (2) to (1);
                \draw (4) to (3);
                \draw (2) to (3);
                \draw (5) to (4);
                \draw (6) to (5);
                \end{tikzpicture}
                \caption{An induced copy of $P_5$ given that $N_\varnothing$ is
                    not anticomplete to $N_X$.}
                \label{fig:pre_3}
            \end{figure}
    \end{enumerate}
\end{proof}

From \Cref{lem:pre_neighborhoods_c5}, it follows that the only exact
neighborhoods $N_X$ of a $5$-cycle $C$ in a $(P_5, \text{house})$-free graph
that can be nonempty are of three types: $X$ consists of three consecutive
vertices of $C$, or $X$ consists of two non-consecutive vertices of $C$, or
$X=V$.  Moreover, the third item in \Cref{lem:pre_neighborhoods_c5} implies that
if $G$ is connected and $N_\varnothing$ is nonempty, then some vertex in
$N_\varnothing$ must be adjacent to some vertex in $N_Z$.

As in the case with chordal graphs, we need to describe relations between the
exact neighborhoods in order to prove the main theorem of the section. Then,
given $X\subseteq Z$, the exact neighborhoods $N_X$ when $X$ is of size 1 will
be called 1-bags. Similarly, $N_X$ is a 2-bag or 3-bag when $X$ is of size 2 or
3, respectively. Moreover, \Cref{lem:pre_neighborhoods_c5} implies that,
essentially, there is only one type of bag for any of the sizes, for example,
the only type of 2-bag that can be nonempty are those that are adjacent to two
nonadjacent vertices of $C$. By that reason, in the following lemmas we only
take into account the bags that can be nonempty, since all others will satisfy
the conclusion of each item vacuously. In particular, notice that any nonempty
3-bag $N_X$ has an associated middle vertex in $X$ to which it is complete.

Additionally, by recalling that the complement of a $5$-cycle is once again a
$5$-cycle (see \Cref{fig:complement}), we observe that if $X \neq Z$, then in
$\overline{G}$ the set $N_X$ is an exact neighborhood of the complement of $C$
of one of the aforementioned types.  This fact implies that, for any two of
these types of neighborhoods, if a relation between them (such as being
anticomplete or complete to each other) is forced due to a $P_5$ or a house,
then the complementary relation must hold for the complement types of
neighborhoods. For example, as we will see in the next lemma, if $N_X$ and $N_Y$
are complete to each other due to a $P_5$, then $N_{Z \setminus X}$ and $N_{Z
\setminus Y}$ are anticomplete to each other, since otherwise one would have a
house, recalling that the complement of $P_5$ is the house.

\begin{figure}[htb!]
    \centering
    \begin{tikzpicture}[every node/.style={vertex}, node
            distance=1cm, on grid, edge]
        \begin{scope}[xshift=-2cm]
            \foreach \i in {0,...,4}
                \node (\i)  at ({(360/5)*\i + 90}:1){};

            \node (5) at (0,0) {};
                
            \foreach \i in {0,...,4}
                \draw let \n1={int(mod(\i+1,5))} in (\i) to (\n1);
            
            \draw (4) to (5);
            \draw (0) to (5);
            \draw (1) to (5);
        \end{scope}

        \begin{scope}[xshift=2cm]
            \foreach \i in {0,...,4}
                \node (\i)  at ({(360/5)*\i + 90}:1){};

            \node (5) at (0,0) {};
                
            \foreach \i in {0,...,4}
                \draw let \n1={int(mod(\i+2,5))} in (\i) to (\n1);
            
            \draw (2) to (5);
            \draw (3) to (5);
        \end{scope}
    \end{tikzpicture}
    \caption{On the left, a vertex of a neighborhood corresponding to three
    consecutive vertices of $C$ is shown. By taking the complement, that vertex
    becomes one that is adjacent to two non-consecutive vertices of
    $\overline{C}$.}
    \label{fig:complement}
\end{figure}

\begin{lemma}
    \label{lem:neighborhoods_c5_complete}
    Let $G$ be a $\p{P_5,\text{house}}$-free graph with an induced copy of
    $C_5$. Let $X$ and $Y$ be different subsets of $Z$.
    \begin{lemenum}
        \item If $N_X$ and $N_Y$ are 2-bags and $X\cap Y$ is empty, then $N_X$
            is complete to $N_Y$. \label{lem:c5_comp_1}
        \item If $N_X$ and $N_Y$ are 3-bags and $|X\cap Y|=2$, then $N_X$ is
            complete to $N_Y$.\label{lem:c5_comp_2}
        \item If $N_X$ is a 3-bag and $N_Y$ is a 2-bag complete to the middle
            vertex of $X$, then $N_X$ is complete to $N_Y$.
            \label{lem:c5_comp_3}
    \end{lemenum}
\end{lemma}

\begin{proof}
    \leavevmode \par
    \begin{enumerate}[label=(\roman*)]
        \item Suppose without loss of generality that $X=\l{a,c}$ and
            $Y=\l{b,d}$. If $x\in N_X$ is nonadjacent to $y\in N_Y$, then
            $\p{a,x,c,d,y}$ is an induced copy of $P_5$. Then, $N_X$ is complete
            to $N_Y$ (see \Cref{fig:complete_1}).

            \begin{figure}[htb!]
                \centering
                \begin{tikzpicture}[every node/.style={vertex}, node
                    distance=1cm, on grid, edge, scale=1.5]
                    \foreach \i in {1,2,4}
                    \node (\i)  at ({(360/5)*\i + 90}:1){};

                    \node[lightgray] (0) at (90:1) {};
                    \node[lightgray] (3) at (306:1) {};

                    \node (5) at (90:0.5) {};
                    \node (6) at (162:0.5) {};
                    
                    \draw[lightgray] (0) to (1);
                    \draw[lightgray] (2) to (3);
                    \draw[lightgray] (3) to (4);
                    \draw[lightgray] (4) to (0);
                    \draw[lightgray] (0) to (6);
                    \draw (2) to (1);
                    \draw (4) to (5);
                    \draw (1) to (5);
                    \draw (2) to (6);
                \end{tikzpicture}
                \caption{An induced copy of $P_5$ given that $N_X$ and $N_Y$ are not complete to each other.}
                \label{fig:complete_1}
            \end{figure}
            
        \item Suppose without loss of generality that $X=\l{a,b,c}$ and
            $Y=\l{b,c,d}$. If $x\in N_X$ is nonadjacent to $y\in N_Y$, then
            $\p{x,b,y,d,e}$ is an induced copy of $P_5$. Thus, $N_X$ is complete
            to $N_Y$ (see \Cref{fig:complete_2}).

            \begin{figure}[htb!]
                \centering
                \begin{tikzpicture}[every node/.style={vertex}, node
                    distance=1cm, on grid, edge, scale=1.5]
                    \foreach \i in {0,2,3}
                    \node (\i)  at ({(360/5)*\i + 90}:1){};

                    \node[lightgray] (1) at (162:1) {};
                    \node[lightgray] (4) at (18:1) {};
                    
                    \node (5) at (90:0.5) {};
                    \node (6) at (162:0.5) {};
                    
                    \draw[lightgray] (0) to (1);
                    \draw[lightgray] (2) to (1);
                    \draw (2) to (3);
                    \draw[lightgray] (3) to (4);
                    \draw[lightgray] (4) to (0);
                    \draw[lightgray] (4) to (5);
                    \draw[lightgray] (1) to (5);
                    \draw (0) to (6);
                    \draw (2) to (6);
                    \draw (0) to (5);
                    \draw[lightgray] (6) to (1);
                \end{tikzpicture}
                \caption{An induced copy of $P_5$ given that $N_X$ and $N_Y$ are not complete to each other.}
                \label{fig:complete_2}
            \end{figure}
    
        \item Suppose without loss of generality that $X=\l{a,b,c}$ and
            $Y=\l{b,d}$. If $x\in N_X$ is nonadjacent to $y\in N_Y$, then
            $\p{x,b,y,d,e}$ is an induced copy of $P_5$. Thus, $N_X$ is complete
            to $N_Y$ (see \Cref{fig:complete_3}).

            \begin{figure}[htb!]
                \centering
                \begin{tikzpicture}[every node/.style={vertex}, node
                    distance=1cm, on grid, edge, scale=1.5]
                    \foreach \i in {0,2,3}
                    \node (\i)  at ({(360/5)*\i + 90}:1){};

                    \node[lightgray] (1) at (162:1) {};
                    \node[lightgray] (4) at (18:1) {};
                    
                    \node (5) at (90:0.5) {};
                    \node (6) at (162:0.5) {};
                    
                    \draw[lightgray] (0) to (1);
                    \draw[lightgray] (2) to (1);
                    \draw (2) to (3);
                    \draw[lightgray] (3) to (4);
                    \draw[lightgray] (4) to (0);
                    \draw[lightgray] (4) to (5);
                    \draw[lightgray] (1) to (5);
                    \draw (0) to (6);
                    \draw (2) to (6);
                    \draw (0) to (5);
                \end{tikzpicture}
                \caption{An induced copy of $P_5$ given that $N_X$ and $N_Y$ are not complete to each other.}
                \label{fig:complete_3}
            \end{figure}
    \end{enumerate}
\end{proof}

The following result states relations of anticompleteness.

\begin{lemma}
    \label{lem:neighborhoods_c5_anti}
    Let $G$ be a $\p{P_5,\text{house}}$-free graph with an induced copy of
    $C_5$. Let $X$ and $Y$ be different subsets of $Z$.
    \begin{lemenum}
        \item If $N_X$ and $N_Y$ are both 3-bags with $|X\cap Y|=1$, then $N_X$
            is anticomplete to $N_Y$. \label{lem:c5_anticomp_1}
        \item If $N_X$ and $N_Y$ are both 2-bags with $|X\cap Y|=1$, then $N_X$
            is anticomplete to $N_Y$. \label{lem:c5_anticomp_2}
        \item If $N_X$ is a 3-bag and $N_Y$ is a 2-bag, where $Y$ intersects $X$
            in just one vertex different from its middle vertex, then $N_X$ is
            anticomplete to $N_Y$. \label{lem:c5_anticomp_3}
        \item If $G$ is $H_{11}$-free, $N_X$ is a 3-bag and $N_Y$ is a 2-bag
            with $|X\cap Y|=2$, then $N_X$ is anticomplete to $N_Y$.
            \label{lem:c5_anticomp_4}
    \end{lemenum}
\end{lemma}
    
\begin{proof}
    \leavevmode\par
    \begin{enumerate}[label=(\roman*)]
        \item Notice $X$ and $Y$ are such that $\overline{X}$ and $\overline{Y}$
            are as in \Cref{lem:c5_comp_1}. Then, since $N_{\overline{X}}$ and
            $N_{\overline{Y}}$ are complete to each other due to an induced copy
            of $P_5$, it follows that $N_X$ is anticomplete to $N_Y$ due to an
            induced copy of the house (see \Cref{fig:anticomplete_1}).

            \begin{figure}[htb!]
                \centering
                \begin{tikzpicture}[every node/.style={vertex}, node
                    distance=1cm, on grid, edge, scale=1.5]
                    \foreach \i in {1,3,4}
                        \node (\i)  at ({(360/5)*\i + 90}:1){};

                    \node[lightgray] (0) at (90:1) {};
                    \node[lightgray] (2) at (234:1) {};
                    
                    \node (5) at (90:0.5) {};
                    \node (6) at (234:0.5) {};
                    
                    \draw[lightgray] (0) to (1);
                    \draw[lightgray] (2) to (1);
                    \draw[lightgray] (2) to (3);
                    \draw (3) to (4);
                    \draw[lightgray] (4) to (0);
                    \draw (4) to (5);
                    \draw (1) to (5);
                    \draw (1) to (6);
                    \draw[lightgray] (2) to (6);
                    \draw (6) to (3);
                    \draw[lightgray] (0) to (5);
                    \draw (6) to (5);
                \end{tikzpicture}
                \caption{An induced copy of the house given that $N_X$ and $N_Y$
                    are not anticomplete to each other.}
                \label{fig:anticomplete_1}
            \end{figure}
        
        \item Notice $X$ and $Y$ are such that $\overline{X}$ and $\overline{Y}$
            are as in \Cref{lem:c5_comp_2}. Then, since $N_{\overline{X}}$ and
            $N_{\overline{Y}}$ are complete to each other due to an induced copy
            of $P_5$, it follows that $N_X$ is anticomplete to $N_Y$ due to an
            induced copy of the house (see \Cref{fig:anticomplete_2}).

            \begin{figure}[htb!]
                \centering
                \begin{tikzpicture}[every node/.style={vertex}, node
                    distance=1cm, on grid, edge, scale=1.5]
                    \foreach \i in {1,3,4}
                        \node (\i)  at ({(360/5)*\i + 90}:1){};

                    \node[lightgray] (0) at (90:1) {};
                    \node[lightgray] (2) at (234:1) {};
                    
                    \node (5) at (90:0.5) {};
                    \node (6) at (234:0.5) {};
                    
                    \draw[lightgray] (0) to (1);
                    \draw[lightgray] (2) to (1);
                    \draw[lightgray] (2) to (3);
                    \draw (3) to (4);
                    \draw[lightgray] (4) to (0);
                    \draw (4) to (5);
                    \draw (1) to (5);
                    \draw (1) to (6);
                    \draw (6) to (3);
                    \draw (6) to (5);
                \end{tikzpicture}
                \caption{An induced copy of the house given that $N_X$ and $N_Y$
                    are not anticomplete to each other.}
                \label{fig:anticomplete_2}
            \end{figure}
        
        \item Notice $X$ and $Y$ are such that $\overline{X}$ and $\overline{Y}$
            are as in \Cref{lem:c5_comp_3}. Then, since $N_{\overline{X}}$ and
            $N_{\overline{Y}}$ are complete to each other due to an induced copy
            of $P_5$, it follows that $N_X$ is anticomplete to $N_Y$ due to an
            induced copy of the house (see \Cref{fig:anticomplete_3}).

            \begin{figure}[htb!]
                \centering
                \begin{tikzpicture}[every node/.style={vertex}, node
                    distance=1cm, on grid, edge, scale=1.5]
                    \foreach \i in {1,3,4}
                        \node (\i)  at ({(360/5)*\i + 90}:1){};

                    \node[lightgray] (0) at (90:1) {};
                    \node[lightgray] (2) at (234:1) {};
                    
                    \node (5) at (90:0.5) {};
                    \node (6) at (234:0.5) {};
                    
                    \draw[lightgray] (0) to (1);
                    \draw[lightgray] (2) to (1);
                    \draw[lightgray] (2) to (3);
                    \draw (3) to (4);
                    \draw[lightgray] (4) to (0);
                    \draw (4) to (5);
                    \draw (1) to (5);
                    \draw (1) to (6);
                    \draw (6) to (3);
                    \draw (6) to (5);
                    \draw[lightgray] (0) to (5);
                \end{tikzpicture}
                \caption{An induced copy of the house given that $N_X$ and $N_Y$
                    are not anticomplete to each other.}
                \label{fig:anticomplete_3}
            \end{figure}
        
        \item Suppose without loss of generality that $X=\l{a,b,c}$, so
            $Y=\l{a,c}$. If $x\in X$ is adjacent to $y\in Y$, notice
            $\l{a,b,c,x,y}$ induces a copy of $H_{11}$. Hence, $N_X$ is
            anticomplete to $N_Y$ (see \Cref{fig:anticomplete_4}).

            \begin{figure}[htb!]
                \centering
                \begin{tikzpicture}[every node/.style={vertex}, node
                    distance=1cm, on grid, edge, scale=1.5]
                    \foreach \i in {0,1,4}
                        \node (\i)  at ({(360/5)*\i + 90}:1){};

                    \node[lightgray] (3) at (306:1) {};
                    \node[lightgray] (2) at (234:1) {};
                    
                    \node (5) at (90:0.5) {};
                    \node (6) {};
                    
                    \draw (0) to (1);
                    \draw[lightgray] (2) to (1);
                    \draw[lightgray] (2) to (3);
                    \draw[lightgray] (3) to (4);
                    \draw (4) to (0);
                    \draw (4) to (5);
                    \draw (1) to (5);
                    \draw (1) to (6);
                    \draw (6) to (4);
                    \draw (6) to (5);
                    \draw (0) to (5);
                \end{tikzpicture}
                \caption{An induced copy of $H_{11}$ given that $N_X$ and $N_Y$
                    are not anticomplete to each other.}
                \label{fig:anticomplete_4}
            \end{figure}
    \end{enumerate}
\end{proof}

The final lemma addresses the structure of these exact neighborhoods, already
considering that $G$ does not contain the aforementioned minimal obstructions.

\begin{lemma}
    \label{lem:neighborhoods_c5_types}
    Let $G$ be a $\p{P_5,\text{house}}$-free graph with an induced copy of $C_5$
    and free of the graphs depicted in \Cref{fig:p5-house-obstruction-with-c5}.
    Let $X$ and $Y$ be different subsets of $Z$.
    \begin{lemenum}
        \item If $N_X$ is a 3-bag, then $N_X$ is a clique. \label{lem:c5_type_1}
        \item If $N_X$ is a 2-bag, then $G[N_X]$ is a cluster.
            \label{lem:c5_type_2}
        \item If $N_X$ and $N_Y$ are as in \Cref{lem:c5_comp_3} and $N_X$ is
            nonempty, then $N_Y$ is stable. \label{lem:c5_type_3}
        \item If $N_X$ and $N_Y$ are as in \Cref{lem:c5_comp_1}, then at most
            one of $N_X$ and $N_Y$ is not stable. \label{lem:c5_type_4}
    \end{lemenum}
\end{lemma}

\begin{proof}
    \leavevmode \par
    \begin{enumerate}[label=(\roman*)]
        \item Suppose without loss of generality that $X=\l{a,b,c}$. Let $x$ and
            $x'$ be nonadjacent vertices in $N_X$. Then, $\l{x,a,x',c,b}$
            induces a copy of $H_{11}$. Therefore, $N_X$ is a clique (see
            \Cref{fig:types_1}).
            
            \begin{figure}[htb!]
                \centering
                \begin{tikzpicture}[every node/.style={vertex}, node
                    distance=1cm, on grid, edge, scale=1.5]
                    \foreach \i in {0,1,4}
                        \node (\i)  at ({(360/5)*\i + 90}:1){};

                    \node[lightgray] (3) at (306:1) {};
                    \node[lightgray] (2) at (234:1) {};
                    
                    \node (5) at ($(0) + (-105:1.25)$) {};
                    \node (6) at ($(0) + (-75:1.25)$){};
                    
                    \draw (0) to (1);
                    \draw[lightgray] (2) to (1);
                    \draw[lightgray] (2) to (3);
                    \draw[lightgray] (3) to (4);
                    \draw (4) to (0);
                    \draw (4) to (5);
                    \draw (1) to (5);
                    \draw (1) to (6);
                    \draw (6) to (4);
                    \draw (6) to (0);
                    \draw (0) to (5);
                \end{tikzpicture}
                \caption{An induced copy of $H_{11}$ given that $N_X$ is not a
                    clique.}
                \label{fig:types_1}
            \end{figure}

        \item Suppose without loss of generality that $X=\l{a,c}$ and suppose
            $G[N_X]$ is not a cluster. Then, there exists $b'$, $b''$ and $b'''$
            such that they induce a copy of $P_3$. Then, $\l{a,b',b'',b''',c}$
            induces a copy of $H_{11}$. Therefore, $G[N_X]$ is a cluster (see
            \Cref{fig:types_2}).

            \begin{figure}[htb!]
                \centering
                \begin{tikzpicture}[every node/.style={vertex}, node
                    distance=1cm, on grid, edge, scale=1.5]
                    \foreach \i in {1,4}
                        \node (\i)  at ({(360/5)*\i + 90}:1){};

                    \node[lightgray] (0) at (90:1) {};
                    \node[lightgray] (3) at (306:1) {};
                    \node[lightgray] (2) at (234:1) {};
                    
                    \node (5) at ($(0) + (0,-0.5)$) {};
                    \node (6) at ($(5) + (0,-0.5)$) {};
                    \node (7) at ($(6) + (0,-0.5)$) {};
                    
                    \draw[lightgray] (0) to (1);
                    \draw[lightgray] (2) to (1);
                    \draw[lightgray] (2) to (3);
                    \draw[lightgray] (3) to (4);
                    \draw[lightgray] (4) to (0);
                    \draw (4) to (5);
                    \draw (1) to (5);
                    \draw (1) to (6);
                    \draw (6) to (4);
                    \draw (6) to (5);
                    \draw (7) to (6);
                    \draw (7) to (1);
                    \draw (7) to (4);
                \end{tikzpicture}
                \caption{An induced copy of $H_{11}$ given that $G[N_X]$ is not
                a cluster.}
                \label{fig:types_2}
            \end{figure}

        \item Suppose without loss of generality that $X=\l{a,b,c}$,
            $Y=\l{b,d}$, and let $b'\in N_X$. If $c'$ and $c''$ in $N_Y$ are
            adjacent, by \Cref{lem:c5_comp_3}, it follows that $b'$ is adjacent
            to both $c'$ and $c''$. Then, $\l{b,b',c,c',c'',d}$ induces a copy
            of $H_{13}$. Hence, $N_Y$ is stable (see \Cref{fig:types_3}).

            \begin{figure}[htb!]
                \centering
                \begin{tikzpicture}[every node/.style={vertex}, node
                    distance=1cm, on grid, edge, scale=1.5]
                    \foreach \i in {0,1,2}
                        \node (\i)  at ({(360/5)*\i + 90}:1){};

                    \node[lightgray] (3) at (306:1) {};
                    \node[lightgray] (4) at (18:1) {};
                    
                    \node (5) at (90:0.5) {};
                    \node (6) at (162:0.5) {};
                    \node (7) at (18:0.5) {};
                    
                    \draw (0) to (1);
                    \draw (2) to (1);
                    \draw[lightgray] (2) to (3);
                    \draw[lightgray] (3) to (4);
                    \draw[lightgray] (4) to (0);
                    \draw[lightgray] (4) to (5);
                    \draw (1) to (5);
                    \draw (0) to (6);
                    \draw (2) to (6);
                    \draw (0) to (5);
                    \draw (5) to (6);
                    \draw (5) to (7);
                    \draw (7) to (6);
                    \draw (7) to (0);
                    \draw (7) to (2);
                \end{tikzpicture}
                \caption{An induced copy of $H_{13}$ given that $N_X$ is
                    nonempty and $N_Y$ is not stable.}
                \label{fig:types_3}
            \end{figure}

        \item Suppose without loss of generality that $X=\l{a,c}$ and
            $Y=\l{b,d}$. If $b'$ and $b''$ are adjacent vertices in $N_X$, and
            $c'$ and $c''$ are adjacent vertices in $N_Y$, then from
            \Cref{lem:c5_comp_1} follows that $\l{b,b',b'',c,c',c''}$ induces a
            copy of $H_{13}$. Thus, one of $N_X$ and $N_Y$ must be stable (see
            \Cref{fig:types_4}).

            \begin{figure}[htb!]
                \centering
                \begin{tikzpicture}[every node/.style={vertex}, node
                    distance=1cm, on grid, edge, scale=1.5]
                    \foreach \i in {0,1}
                        \node (\i)  at ({(360/5)*\i + 90}:1){};

                    \node[lightgray] (2) at (234:1) {};
                    \node[lightgray] (3) at (306:1) {};
                    \node[lightgray] (4) at (18:1) {};
                    
                    \node (5) at ($(1) + (86:0.5)$){};
                    \node (7) at ($(0) + (166:0.5)$) {};
                    \node (6) at ($(7) + (126:0.5)$) {};
                    \node (8) at ($(5) + (126:0.5)$) {};
                    
                    \draw[lightgray] (2) to (6);
                    \draw[lightgray] (7) to (2);
                    \draw[lightgray] (2) to (1);
                    \draw[lightgray] (2) to (3);
                    \draw[lightgray] (3) to (4);
                    \draw[lightgray] (4) to (0);
                    \draw[lightgray] (4) to (5);
                    \draw[lightgray] (8) to (4);
                    \draw (1) to (5);
                    \draw (0) to (1);
                    \draw (0) to (6);
                    \draw (5) to (6);
                    \draw (5) to (7);
                    \draw (7) to (6);
                    \draw (7) to (0);
                    \draw (8) to (1);
                    \draw (8) to (5);
                    \draw (8) to (6);
                    \draw (8) to (7);
                \end{tikzpicture}
                \caption{An induced copy of $H_{13}$ given that none of $N_X$
                    and $N_Y$ are stable.}
                \label{fig:types_4}
            \end{figure}
    \end{enumerate}
\end{proof}

We proceed to prove that the claimed minimal obstructions are indeed all the
graphs that prevent $G$ from being monopolar.

\begin{theorem}
    \label{thm:monopolar_c5}
    A connected $(P_5,\text{house})$-free graph $G$ with an induced copy of
    $C_5$ is monopolar if and only if it does not contain any of the graphs
    depicted in
    \Cref{fig:chordal-p5-free-obstructions,fig:p5-house-obstruction-with-c5} as
    an induced subgraph.
\end{theorem}

\begin{proof}
    Since monopolarity is hereditary, it follows from
    \Cref{prop:minimal_obs_c5,prop:minimal_obs_chordal} that $G$ does not
    contain the mentioned graphs as induced subgraphs.
    
    Due to \Cref{lem:pre_c5_1,lem:pre_c5_2}, the only exact neighborhoods of $C$
    which can be nonempty are the neighborhoods of two non-consecutive vertices,
    and the neighborhoods of three consecutive vertices.   Since $G$ is
    $H_8$-free, $N_Z$ is empty, but by \Cref{lem:pre_c5_3} and the connectedness
    of $G$, we conclude that $N_\varnothing$ is also empty. Hence, the only
    possibly nonempty sets are
    \[
        N_{ac}, N_{ad}, N_{bd}, N_{be}, N_{ce}, N_{abc}, N_{bcd}, N_{cde},
        N_{ade} \text{ and } N_{abe}.
    \]
    
    Suppose that there are four different nonempty $3$-bags. Notice that three
    of them must be complete to the same single vertex, so suppose without loss
    of generality that $W=\l{a,b,c}$, $X=\l{b,c,d}$ and $Y=\l{c,d,e}$. Let $w\in
    N_W$, $x\in N_X$ and $y\in N_Y$. On the one hand, \Cref{lem:c5_comp_2}
    implies that $x$ is adjacent to both $w$ and $y$. On the other hand,
    \Cref{lem:c5_anticomp_1} ensures that $w$ and $y$ are not adjacent.
    Therefore, $\l{w,x,y,b,c,d}$ induces a copy of $H_1$ (see
    \Cref{fig:3_consec_bags}). Thus, it cannot be that more than three 3-bags
    are nonempty. We proceed by first considering the case in which three of
    these are nonempty.

    \begin{figure}[htb!]
        \centering
        \begin{tikzpicture}[every node/.style={vertex}, node
            distance=1cm, on grid, edge, scale=1.5]
            \foreach \i in {0,1,2}
            \node (\i)  at ({(360/5)*\i + 90}:1){};

            \node[lightgray] (3) at (306:1) {};
            \node[lightgray] (4) at (18:1) {};
            
            \node (5) at (90:0.5) {};
            \node (6) at (162:0.5) {};
            \node (7) at (234:0.5) {};
            
            \draw[lightgray] (2) to (3);
            \draw[lightgray] (7) to (3);
            \draw[lightgray] (3) to (4);
            \draw[lightgray] (4) to (0);
            \draw[lightgray] (4) to (5);
            \draw (1) to (5);
            \draw (0) to (1);
            \draw (2) to (1);
            \draw (6) to (1);
            \draw (0) to (6);
            \draw (2) to (6);
            \draw (0) to (5);
            \draw (7) to (1);
            \draw (7) to (2);
            \draw (7) to (6);
            \draw (5) to (6);
        \end{tikzpicture}
        \caption{An induced copy of $H_1$ given that four 3-bags are nonempty.}
        \label{fig:3_consec_bags}
    \end{figure}
    
    \paragraph{Case 1: Three $3$-bags are nonempty.} The aforementioned argument
    implies that no vertex of $C$ is adjacent to all the $3$-bags. Hence,
    suppose without loss of generality that these are $N_{abc}$, $N_{bcd}$ and
    $N_{ade}$. Take $b'\in N_{abc}$, $c'\in N_{bcd}$ and $e'\in N_{ade}$. Once
    more, \Cref{lem:c5_comp_2} implies that $b'$ is adjacent to $c'$, while
    \Cref{lem:c5_anticomp_1} implies that $e'$ is nonadjacent to both $a'$ and
    $b'$. If $b''\in N_{ac}$, \Cref{lem:c5_comp_3} implies that $b''$ is
    adjacent to $c'$. Additionally, \Cref{lem:c5_anticomp_3,lem:c5_anticomp_4}
    imply that $b''$ is nonadjacent to both $b'$ and $e'$. Then,
    $\l{a,b,b',c,c',b''}$ induces a copy of $H_{13}$ (see
    \Cref{fig:2_bag_nonempty}). It follows that $N_{ac}$ is empty, and
    analogously $N_{bd}$ is empty.

    \begin{figure}[htb!]
        \centering
        \begin{tikzpicture}[every node/.style={vertex}, node
            distance=1cm, on grid, edge, scale=1.5]
            \foreach \i in {0,1,4}
                \node (\i)  at ({(360/5)*\i + 90}:1){};

            \node[lightgray] (3) at (306:1) {};
            \node[lightgray] (2) at (234:1) {};
            
            \node (5) at (90:0.5) {};
            \node (6) at (162:0.5) {};
            \node[lightgray] (7) at (306:0.5) {};
            \node (8) at (-0.25,-0.25) {};
            
            \draw[lightgray] (2) to (3);
            \draw[lightgray] (7) to (3);
            \draw[lightgray] (3) to (4);
            \draw[lightgray] (2) to (1);
            \draw[lightgray] (2) to (6);
            \draw[lightgray] (7) to (4);
            \draw[lightgray] (7) to (2);
            \draw (4) to (0);
            \draw (4) to (5);
            \draw (1) to (5);
            \draw (0) to (1);
            \draw (6) to (1);
            \draw (0) to (6);
            \draw (0) to (5);
            \draw (5) to (6);
            \draw (8) to (1);
            \draw (8) to (4);
            \draw (8) to (6);
        \end{tikzpicture}
        \caption{An induced copy of $H_{13}$ given that $N_{ac}$ is nonempty.}
        \label{fig:2_bag_nonempty}
    \end{figure}

    By \Cref{lem:c5_type_3}, $N_{be}$ and $N_{ce}$ are stable sets. Moreover, by
    \Cref{lem:c5_anticomp_2}, $N_{be}$ is anticomplete to $N_{ce}$. Observe that
    \[
        N_{be} \cup N_{ce} \cup \l{a,d}
    \]
    is a stable set.

    Lastly, by \Cref{lem:c5_type_2}, $G[N_{ad}]$ is a cluster such that, by
    \Cref{lem:c5_anticomp_3,lem:c5_anticomp_4}, is anticomplete to $N_{abc}$,
    $N_{bcd}$ and $N_{ade}$. Meanwhile, $N_{abc}$ is complete to $N_{bcd}$ by
    \Cref{lem:c5_comp_2}, and \Cref{lem:c5_anticomp_1} ensures that $N_{ade}$ is
    anticomplete to both $N_{abc}$ and $N_{bcd}$. Also, all 3-bags are cliques
    due to \Cref{lem:c5_type_1}, so 
    \[
        N_{abc} \cup N_{bcd} \cup N_{ade} \cup N_{ad} \cup \l{b,c,e}
    \]
    induces a cluster. This proves that
    \[
        \p{N_{be} \cup N_{ce} \cup \l{a,d} , N_{abc} \cup N_{bcd} \cup N_{ade} 
        \cup N_{ad} \cup \l{b,c,e}}
    \]
    is a monopolar partition of $G$. For completion of the relations between
    sets, we have that $N_{be}$ is complete to $N_{abc}$ and $N_{ade}$ by
    \Cref{lem:c5_comp_3}, while $N_{be}$ is complete to $N_{ad}$ by
    \Cref{lem:c5_comp_1}. Similarly, for $N_{ce}$ we have that it is complete to
    $N_{bcd}$, $N_{ade}$ and $N_{ad}$. Therefore, in
    \Cref{fig:monopolar_partition_1} we show the monopolar partition of $G$,
    along with its complete structure.

    \begin{figure}[htb!]
        \centering
        \begin{tikzpicture}[every node/.style={vertex}, node
            distance=1cm, on grid, edge, scale=1.5]
            \foreach \i in {0,1,4}
            \node (\i)  at ({(360/5)*\i + 90}:1){};

            \node (3) at (306:1) {};
            \node (2) at (234:1) {};
            
            \node (5) at (90:0.5) {};
            \node (6) at (162:0.5) {};
            \node (7) at (306:0.75) {};
            \node (8) at (306:0.5) {};
            \node (9) at (18:0.5) {};
            \node (10) at (234:0.5) {};
            
            \draw (2) to (3);
            \draw[shorten <=1pt, shorten >=1pt] (7) to (3);
            \draw (3) to (4);
            \draw (2) to (1);
            \draw (2) to (6);
            \draw (7) to (4);
            \draw (7) to (2);
            \draw (4) to (0);
            \draw (4) to (5);
            \draw (1) to (5);
            \draw (0) to (1);
            \draw (6) to (1);
            \draw (0) to (6);
            \draw (0) to (5);
            \draw (5) to (6);
            \draw (8) to (2);
            \draw (8) to (4);
            \draw (9) to (0);
            \draw (9) to (3);
            \draw (10) to (3);
            \draw (10) to (1);
            \draw (9) to (5);
            \draw (9) to (7);
            \draw (9) to (8);
            \draw (10) to (7);
            \draw (10) to (8);
            \draw (10) to (6);

            \begin{scope}[on background layer]
                \draw[fill=red!20, draw=none] \convexpath{1,0,5,6}{3.5pt};
                \draw[fill=red!20, draw=none] \convexpath{7,3}{3.5pt};
                \node [fit = (8), fill=red!20, draw=none, inner sep=0.5pt]
                {}; 
                \draw[fill=blue!20, draw=none] \convexpath{2,10}{3.5pt};
                \draw[fill=blue!20, draw=none] \convexpath{4,9}{3.5pt};
            \end{scope}
        \end{tikzpicture}
        \caption{A monopolar partition of $G$.}
        \label{fig:monopolar_partition_1}
    \end{figure}

    \paragraph{Case 2: Two $3$-bags are not empty.} If the nonempty 3-bags are
    such that two vertices of $C$ are complete to both, then the partition of
    the previous case still is a monopolar partition for $G$, noting that any of
    the obstructions mentioned in that case did not use the fact that $N_{ade}$
    was nonempty. Then, suppose without loss of generality that $N_{abc}$ and
    $N_{cde}$ are the nonempty 3-bags.

    Let $b'\in N_{abc}$ and $d'\in N_{cde}$. Consider as a subcase that one of
    $N_{ac}$ and $N_{ce}$ is nonempty. If both were to be nonempty, take $b''\in
    N_{ac}$ and $d''\in N_{ce}$. By \Cref{lem:c5_anticomp_2}, $b''$ is not
    adjacent to $d''$, while \Cref{lem:c5_anticomp_4} ensures that $b''$ is not
    adjacent to $b'$. Additionally, \Cref{lem:c5_anticomp_3} ensures that $d''$
    is not adjacent to $b'$. Thus, $\l{a,b,c,d,e,b',b'',d''}$ induces a copy of
    $H_9$ (see \Cref{fig:two_2_bags_nonempty}).

    \begin{figure}[htb!]
        \centering
        \begin{tikzpicture}[every node/.style={vertex}, node
            distance=1cm, on grid, edge]
            \begin{scope}[scale=1.5]
                \foreach \i in {0,...,4}
                \node (\i)  at ({(360/5)*\i + 90}:1){};
                
                \node (5) at (90:0.65) {};
                \node[lightgray] (6) at (234:0.65) {};
                \node (7) at (90:0.4) {};
                \node (8) at (234:0.4) {};
                
                \draw (2) to (3);
                \draw (3) to (4);
                \draw (2) to (1);
                \draw (7) to (4);
                \draw (7) to (1);
                \draw (4) to (0);
                \draw (4) to (5);
                \draw (1) to (5);
                \draw (0) to (1);
                \draw (0) to (5);
                \draw (8) to (3);
                \draw (8) to (1);
                \draw[lightgray] (6) to (1);
                \draw[lightgray] (3) to (6);
                \draw[lightgray] (2) to (6);
            \end{scope}
        \end{tikzpicture}
        \caption{An induced copy of $H_9$ given that both $N_{ac}$ and $N_{ce}$ are nonempty.}
        \label{fig:two_2_bags_nonempty}
    \end{figure}

    As a consequence, suppose without loss of generality that $N_{ce}$ is empty.
    Analogously, it follows that $N_{ad}$ is empty.

    By \Cref{lem:c5_type_3}, $N_{bd}$ and $N_{be}$ are stable sets, while
    \Cref{lem:c5_anticomp_2} implies that they are anticomplete to each other.
    
    On the one hand,
    \[
        N_{bd} \cup N_{be} \cup \l{a,c}
    \]
    is a stable set. On the other hand, \Cref{lem:c5_type_2} implies that
    $G[N_{ac}]$ is a cluster. Meanwhile,
    \Cref{lem:c5_anticomp_4,lem:c5_anticomp_3} imply that $N_{ac}$ is
    anticomplete to $N_{abc}$ and $N_{cde}$ respectively, and
    \Cref{lem:c5_anticomp_1} implies that $N_{abc}$ is anticomplete to
    $N_{cde}$. Additionally, $N_{abc}$ and $N_{cde}$ are cliques due to
    \Cref{lem:c5_type_1}. Therefore, 
    \[
        N_{abc}\cup N_{cde} \cup N_{ac} \cup \l{b,d,e}
    \]
    induces a cluster, so
    \[
        \p{N_{bd} \cup N_{be} \cup \l{a,c} , N_{abc}\cup N_{cde} \cup N_{ac} 
        \cup \l{b,d,e}}
    \]
    is a monopolar partition of $G$. As with the previous case, we can exhaust
    the relations between the existing bags. Notice that \Cref{lem:c5_comp_1}
    implies that $N_{ac}$ is complete to both $N_{bd}$ and $N_{be}$, while
    \Cref{lem:c5_comp_3} implies that $N_{bd}$ is complete to $N_{abc}$,
    $N_{bd}$ is complete to $N_{cde}$, and $N_{be}$ is complete to $N_{abc}$.
    Finally, \Cref{lem:c5_anticomp_3} ensures that $N_{be}$ is anticomplete to
    $N_{cde}$. Hence, in \Cref{fig:monopolar_partition_2} we show the monopolar
    partition of $G$, along with its complete structure.

    \begin{figure}[htb!]
        \centering
        \begin{tikzpicture}[every node/.style={vertex}, node
            distance=1cm, on grid, edge]
            \begin{scope}[scale=1.5]
                \foreach \i in {0,...,4}
                \node (\i)  at ({(360/5)*\i + 90}:1){};
                
                \node (5) at (90:0.65) {};
                \node (6) at (234:0.65) {};
                \node (7) at (90:0.4) {};
                \node (8) at (162:0.65) {};
                \node (9) at (18:0.65) {};
                
                \draw (2) to (3);
                \draw (3) to (4);
                \draw (2) to (1);
                \draw (7) to (4);
                \draw (7) to (1);
                \draw (4) to (0);
                \draw (4) to (5);
                \draw (1) to (5);
                \draw (0) to (1);
                \draw (0) to (5);
                \draw (8) to (2);
                \draw (8) to (0);
                \draw (9) to (0);
                \draw (9) to (3);
                \draw (6) to (1);
                \draw (3) to (6);
                \draw (2) to (6);
                \draw (5) to (8);
                \draw (5) to (9);
                \draw (7) to (8);
                \draw (7) to (9);
                \draw (6) to (8);

                \begin{scope}[on background layer]
                    \draw[fill=blue!20, draw=none] \convexpath{1,8}{3.5pt};
                    \draw[fill=blue!20, draw=none] \convexpath{9,4}{3.5pt};
                    \draw[fill=red!20, draw=none] \convexpath{2,6,3}{3.5pt};
                    \draw[fill=red!20, draw=none] \convexpath{0,5}{3.5pt};
                    \node[fit = (7), draw=none, fill=red!20, inner sep=0.5pt] {};
                \end{scope}
            \end{scope}
        \end{tikzpicture}
        \caption{A monopolar partition of $G$.}
        \label{fig:monopolar_partition_2}
    \end{figure}

    On the contrary, assume both $N_{ac}$ and $N_{ce}$ are empty. Observe that
    one of $N_{be}$ or $N_{ad}$ is empty, otherwise, if $a'\in N_{be}$ and
    $e'\in N_{ad}$, by \Cref{lem:c5_comp_3}, $a'$ is adjacent to $b'$ and $e'$
    is adjacent to $d'$, while \Cref{lem:c5_comp_1} implies that $a'$ is
    adjacent to $e'$. Furthermore, \Cref{lem:c5_anticomp_1} ensures that $b'$ is
    nonadjacent to $d'$. Hence, $\l{a,a',b,b',d,d',e,e'}$ induces a copy of
    $H_{12}$ (see \Cref{fig:two_2_bags_nonempty_2}).

    \begin{figure}[htb!]
        \centering
        \begin{tikzpicture}[every node/.style={vertex}, node
            distance=1cm, on grid, edge]
            \begin{scope}[scale=1.5]
                \foreach \i in {0,2,3,4}
                \node (\i)  at ({(360/5)*\i + 90}:1){};
                
                \node[lightgray] (1) at (162:1){};
                \node (5) at (90:0.65) {};
                \node (6) at (234:0.65) {};
                \node (7) at (18:0.65) {};
                \node (8) at (306:0.65) {};
                
                \draw[lightgray] (1) to (5);
                \draw[lightgray] (0) to (1);
                \draw[lightgray] (2) to (1);
                \draw[lightgray] (6) to (1);
                \draw (3) to (6);
                \draw (2) to (6);
                \draw (2) to (3);
                \draw (3) to (4);
                \draw (7) to (3);
                \draw (7) to (0);
                \draw (4) to (0);
                \draw (4) to (5);
                \draw (0) to (5);
                \draw (8) to (4);
                \draw (8) to (2);
                \draw (6) to (8);
                \draw (7) to (8);
                \draw (7) to (5);
            \end{scope}
        \end{tikzpicture}
        \caption{An induced copy of $H_{12}$ given that both $ N_{be}$ and
            $N_{ad}$ are nonempty.}
        \label{fig:two_2_bags_nonempty_2}
    \end{figure}
    
    Hence, suppose without loss of generality that $N_{ad}$ is empty. Thus, the
    monoplar partition for the previous case works as well for this case (with
    $N_{ac}$ empty).
    
    \paragraph{Case 3: Only one $3$-bag is nonempty.} Suppose without loss of
    generality that $N_{abc}$ is nonempty. The following facts from the previous
    case still hold:

    \begin{itemize}
        \item $N_{bd}$ and $N_{be}$ are stable and anticomplete to each other,
        \item $N_{bd}$ and $N_{be}$ are complete to both $N_{abc}$ and $N_{ac}$,
        \item $N_{ac}$ is anticomplete to $N_{abc}$,
        \item $G[N_{ac}]$ is a cluster and
        \item $N_{abc}$ is a clique.
    \end{itemize}

    If $N_{ac}$ is nonempty, as in the previous case, it follows that $N_{ad}$
    and $N_{ce}$ are empty. Therefore, the last monopolar partition given works
    in this case, taking $N_{cde}$ to be empty. 
    
    If $N_{ac}$ is empty, suppose first that $N_{bd}$ and $N_{be}$ are nonempty.
    Let $c'\in N_{bd}$ and $a'\in N_{be}$. If $d'\in N_{ce}$, then
    \Cref{lem:c5_comp_1} implies that $d'$ is adjacent to $c'$, while
    \Cref{lem:c5_anticomp_2} implies that $d'$ is not adjacent to $a'$. Thus,
    $\l{a,a',b,c,c',d,d',e}$ induces a copy of $H_{10}$ (see
    \Cref{fig:2_bag_nonempty_2}).

    \begin{figure}[htb!]
        \centering
        \begin{tikzpicture}[every node/.style={vertex}, node
            distance=1cm, on grid, edge]
            \begin{scope}[scale=1.5]
                \foreach \i in {0,...,4}
                \node (\i)  at ({(360/5)*\i + 90}:1){};
                
                \node[lightgray] (5) at (90:0.65) {};
                \node (6) at (18:0.65) {};
                \node (7) at (162:0.65) {};
                \node (8) at (234:0.65) {};
                
                \draw[lightgray] (0) to (5);
                \draw[lightgray] (1) to (5);
                \draw[lightgray] (4) to (5);
                \draw (2) to (3);
                \draw (3) to (4);
                \draw (2) to (1);
                \draw (0) to (7);
                \draw (2) to (7);
                \draw (4) to (0);
                \draw (0) to (1);
                \draw (8) to (3);
                \draw (8) to (1);
                \draw (0) to (6);
                \draw (3) to (6);
                \draw (7) to (8);
            \end{scope}
        \end{tikzpicture}
        \caption{An induced copy of $H_{10}$ given that $N_{ce}$ is nonempty.}
        \label{fig:2_bag_nonempty_2}
    \end{figure}
    
    It follows that $N_{ce}$ is empty, and analogously $N_{ad}$ is empty. Once
    more, the previous monopolar partition works as well, taking both $N_{cde}$
    and $N_{ac}$ to be empty.

    Suppose without loss of generality that $N_{bd}$ is empty. If $N_{be}$ is
    nonempty, let $a'\in N_{be}$ and suppose $d'$ and $d''$ are adjacent
    vertices in $N_{ce}$. By \Cref{lem:c5_anticomp_2}, $a'$ is not adjacent to
    both $d'$ and $d''$. Then, $\l{a,a',b,c,d,d',d'',e}$ induces a copy of $H_9$
    (see \Cref{fig:2_bag_not_stable}).

    \begin{figure}[htb!]
        \centering
        \begin{tikzpicture}[every node/.style={vertex}, node
            distance=1cm, on grid, edge]
            \begin{scope}[scale=1.5]
                \foreach \i in {0,...,4}
                \node (\i)  at ({(360/5)*\i + 90}:1){};
                
                \node[lightgray] (5) at (90:0.65) {};
                \node (6) at (18:0.65) {};
                \node (8) at (234:0.65) {};
                \node (7) {};
                
                \draw[lightgray] (0) to (5);
                \draw[lightgray] (1) to (5);
                \draw[lightgray] (4) to (5);
                \draw (2) to (3);
                \draw (3) to (4);
                \draw (2) to (1);
                \draw (1) to (7);
                \draw (3) to (7);
                \draw (4) to (0);
                \draw (0) to (1);
                \draw (8) to (3);
                \draw (8) to (1);
                \draw (0) to (6);
                \draw (3) to (6);
                \draw (7) to (8);
            \end{scope}
        \end{tikzpicture}
        \caption{An induced copy of $H_9$ given that $N_{ce}$ is not stable.}
        \label{fig:2_bag_not_stable}
    \end{figure}
    
    It follows that $N_{ce}$ is stable. From \Cref{lem:c5_anticomp_2} we have
    that $N_{ce}$ is anticomplete to $N_{bd}$, while \Cref{lem:c5_type_3}
    implies that $N_{be}$ is stable, so
    \[
        N_{be} \cup N_{ce} \cup \l{a,d}
    \]
    is a stable set. Observe that, by \Cref{lem:c5_type_2}, $G[N_{ad}]$ is a
    cluster. Also, \Cref{lem:c5_anticomp_3} implies that $N_{abc}$ is
    anticomplete to both $N_{ad}$ and $N_{ce}$. Therefore,
    \[
        N_{abc} \cup N_{ad} \cup \l{b,c,e}
    \]
    is a cluster, so
    \[
        \p{N_{be} \cup N_{ce} \cup \l{a,d}, N_{abc} \cup N_{ad} \cup \l{b,c,e}}
    \]
    is a monopolar partition of $G$. Moreover, \Cref{lem:c5_comp_1} implies that
    $N_{ad}$ is complete to $N_{be}$ and $N_{ce}$, so
    \Cref{fig:monopolar_partition_3} shows a monopolar partition of $G$ along
    with its complete structure.

    \begin{figure}[htb!]
        \centering
        \begin{tikzpicture}[every node/.style={vertex}, node
            distance=1cm, on grid, edge]
            \begin{scope}[scale=1.5]
                \foreach \i in {0,...,4}
                \node (\i)  at ({(360/5)*\i + 90}:1){};
                
                \node (5) at (90:0.65) {};
                \node (7) at (234:0.65) {};
                \node (8) at (306:0.65) {};
                \node (6) at (18:0.65) {};
                
                \draw (2) to (3);
                \draw (3) to (4);
                \draw (2) to (1);
                \draw (7) to (1);
                \draw (7) to (3);
                \draw (4) to (0);
                \draw (4) to (5);
                \draw (1) to (5);
                \draw (0) to (1);
                \draw (0) to (5);
                \draw (8) to (2);
                \draw (8) to (4);
                \draw (6) to (0);
                \draw (6) to (3);
                \draw (3) to (6);
                \draw (5) to (6);
                \draw (7) to (8);
                \draw (6) to (8);

                \begin{scope}[on background layer]
                    \draw[fill=blue!20, draw=none] \convexpath{4,6}{3.5pt};
                    \draw[fill=blue!20, draw=none] \convexpath{2,7}{3.5pt};
                    \draw[fill=red!20, draw=none] \convexpath{1,0,5}{3.5pt};
                    \node[fit = (3), draw=none, fill=red!20, inner sep=0.5pt]
                    {};
                    \node[fit = (8), draw=none, fill=red!20, inner sep=0.5pt] {};
                \end{scope}
            \end{scope}
        \end{tikzpicture}
        \caption{A monopolar partition of $G$.}
        \label{fig:monopolar_partition_3}
    \end{figure}

    Lastly, consider $N_{be}$ to be empty. By \Cref{lem:c5_type_4}, at least one
    between $N_{ad}$ and $N_{ce}$ is a stable set. Suppose without loss of
    generality that $N_{ce}$ is stable. Then, the previous partition works as
    well, taking $N_{be}$ to be empty.

    \paragraph{Case 4: All 3-bags are empty.} If some 2-bag is not stable, say
    $N_{ac}$ has $b'$ and $b''$ adjacent, then taking the cycle
    $\p{a,b',c,d,e}$, one can check that it has a nonempty 3-bag, since $b''$ is
    in $N_{ab'c}$ (see \Cref{fig:2_bag_not_stable_2}). Thus, one can proceed as
    in the previous case. Therefore, consider all $2$-bags to be stable.

    \begin{figure}[htb!]
        \centering
        \begin{tikzpicture}[every node/.style={vertex}, node
            distance=1cm, on grid, edge]
            \begin{scope}[scale=1.5]
                \foreach \i in {1,...,4}
                \node (\i)  at ({(360/5)*\i + 90}:1){};

                \node[lightgray] (0) at (90:1) {};
                
                \node (5) at (90:0.5) {};
                \node (6) {};
                
                \draw[lightgray] (0) to (1);
                \draw[lightgray] (4) to (0);
                \draw (2) to (1);
                \draw (2) to (3);
                \draw (3) to (4);
                \draw (4) to (5);
                \draw (1) to (5);
                \draw (1) to (6);
                \draw (6) to (4);
                \draw (6) to (5);
            \end{scope}
        \end{tikzpicture}
        \caption{An induced copy of $C_5$ with a non-stable 2-bag induces a
            nonempty 3-bag.}
        \label{fig:2_bag_not_stable_2}
    \end{figure}
    
    Suppose at least four 2-bags are nonempty. Without loss of generality,
    suppose $N_{ac}$, $N_{ce}$ and $N_{be}$ are nonempty. Take $b'\in N_{ac}$,
    $d'\in N_{ce}$ and $a'\in N_{be}$. By \Cref{lem:c5_anticomp_2}, $d'$ is
    nonadjacent to both $b'$ and $a'$, and these two are adjacent to each other
    by \Cref{lem:c5_comp_1}. Then, $\l{a,a',b,b',c,d,d',e}$ induces a copy of
    $H_{10}$ (see \Cref{fig:four_2_bags}).

    \begin{figure}[htb!]
        \centering
        \begin{tikzpicture}[every node/.style={vertex}, node
            distance=1cm, on grid, edge]
            \begin{scope}[scale=1.5]
                \foreach \i in {0,...,4}
                \node (\i)  at ({(360/5)*\i + 90}:1){};
                
                \node (5) at (90:0.65) {};
                \node (6) at (234:0.65) {};
                \node (7) at (18:0.65) {};
                
                \draw (0) to (1);
                \draw (4) to (0);
                \draw (2) to (1);
                \draw (2) to (3);
                \draw (3) to (4);
                \draw (4) to (5);
                \draw (1) to (5);
                \draw (1) to (6);
                \draw (3) to (6);
                \draw (0) to (7);
                \draw (3) to (7);
                \draw (5) to (7);
            \end{scope}
        \end{tikzpicture}
        \caption{An induced copy of $H_{10}$ given that at least four $2$-bags
            are nonempty.}
        \label{fig:four_2_bags}
    \end{figure}
    
    This implies that at most three 2-bags are nonempty. If three of them happen
    to be nonempty, a vertex of $C$ must be complete to two of them. Suppose
    without loss of generality that $N_{ac}$ and $N_{ad}$ are such 2-bags. The
    aforementioned argument implies that $N_{bd}$ and $N_{ce}$ are empty. Due to
    \Cref{lem:c5_anticomp_2}, $N_{ac}$ is anticomplete to $N_{ad}$. Moreover,
    since all 2-bags in this case are stable, we obtain
    \[
        \p{N_{ac} \cup N_{ad} \cup \l{b,e} , N_{be} \cup \l{a,c,d}}
    \]
    as a monopolar partition of $G$. Notice that \Cref{lem:c5_comp_1} implies
    that $N_{be}$ is complete to $N_{ac}$ and $N_{ad}$, so in
    \Cref{fig:monopolar_partition_5} the monopolar partition of $G$, along with
    its complete structure, is shown.
    
    \begin{figure}[htb!]
        \centering
        \begin{tikzpicture}[every node/.style={vertex}, node
            distance=1cm, on grid, edge]
            \begin{scope}[scale=1.5]
                \foreach \i in {0,...,4}
                \node (\i)  at ({(360/5)*\i + 90}:1){};
                
                \node (5) at (90:0.65) {};
                \node (6) at (306:0.65) {};
                \node (7) at (18:0.65) {};
                
                \draw (0) to (1);
                \draw (4) to (0);
                \draw (2) to (1);
                \draw (2) to (3);
                \draw (3) to (4);
                \draw (4) to (5);
                \draw (1) to (5);
                \draw (2) to (6);
                \draw (4) to (6);
                \draw (0) to (7);
                \draw (3) to (7);
                \draw (5) to (7);
                \draw (6) to (7);

                \begin{scope}[on background layer]
                    \draw[fill=blue!20, draw=none] \convexpath{0,5}{3.5pt};
                    \draw[fill=blue!20, draw=none] \convexpath{3,6}{3.5pt};
                    \draw[fill=red!20, draw=none] \convexpath{1,2}{3.5pt};
                    \node [fit=(4), fill=red!20, draw=none, inner sep=0.5pt] {}; 
                    \node [fit=(7), fill=red!20, draw=none, inner sep=0.5pt] {}; 
                \end{scope}
            \end{scope}
        \end{tikzpicture}
        \caption{A monopolar partition of $G$.}
        \label{fig:monopolar_partition_5}
    \end{figure}

    If there were less than three nonempty 2-bags, notice the same partition is
    still a monopolar partition of $G$. At last, we conclude that $G$ is
    monopolar.
    
\end{proof}

Observe that we have shown, exhaustively, how the monopolar graphs in this
family look like. Thus, the theorem yields the following corollary.

\begin{corollary}
    A connected monopolar graph $(P_5,\text{house})$-free with an induced copy
    of $C_5$ has the structure of one of the graphs shown in
    \Cref{fig:monopolar_partition_1,fig:monopolar_partition_2,fig:monopolar_partition_3,fig:monopolar_partition_5}.
\end{corollary}